\originalchapter{双陪集与二面体群表示}
\originalsection{Mackey分解}
诱导后再限制到另一个子群 $K$, 陪集块不一定仍被 $K$ 传递地置换. 应把这些块按 $K$ 的轨道分组. $K$ 在 $G/H$ 上的轨道恰由双陪集 $KtH$ 参数化. 双陪集指所有 $kth$ ($k\in K,h\in H$) 组成的集合, 两个这样的集合或者相同或者不交.

\begin{theorem}[Mackey分解]\label{thm:mackey}
取 $K\backslash G/H$ 的代表 $t$. 令 $L_t=K\cap tHt^{-1}$, 在 $W$ 的原空间上定义 $L_t$ 表示 $W_t$, 使 $\ell$ 的作用为 $\rho_W(t^{-1}\ell t)$. 则
\[
 \Res_K^G\Ind_H^G W\cong
 \bigoplus_{t\in K\backslash G/H}\Ind_{L_t}^K W_t.
\]
\end{theorem}
\begin{proof}
依照 $K$ 在左陪集上的轨道, 把 $\Ind_H^G W$ 的陪集直和分成若干 $K$ 不变子空间 $M_t$, 分别由 $kt\otimes w$ 张成. 定义
\[
 \C[K]\otimes_{\C[L_t]} W_t\longrightarrow M_t,\qquad
 k\otimes w\longmapsto kt\otimes w.
\]
若 $\ell\in L_t$, 则 $k\ell t=kt(t^{-1}\ell t)$, 因而 $(k\ell)t\otimes w=kt\otimes\rho_W(t^{-1}\ell t)w$. 这正是源中的平衡关系, 所以映射良定义且交织. 它显然满射. 左陪集 $kL_t$ 与轨道中的陪集 $ktH$ 一一对应: 相同条件都是 $k_2^{-1}k_1\in K\cap tHt^{-1}$. 所以两边维数均为 $[K:L_t]\dim W$, 满射即同构. 所有轨道的直和给出结论.
\end{proof}
旧讲义§4.9把诱导特征标的陪集求和称为Mackey公式; 这里进一步给出双陪集版本, 是编者为计算诱导不可约性补充的结构论证. 两者均来自同一个陪集分块.

\originalsection{诱导不可约性的判据}
\begin{corollary}\label{cor:mackeycriterion}
若 $W$ 是 $H$ 不可约表示, 则 $\Ind_H^G W$ 不可约, 当且仅当对每个非单位双陪集代表 $t$ 有
\[
 \Hom_{H\cap tHt^{-1}}(\Res_{H\cap tHt^{-1}}^H W,W_t)=0.
\]
\end{corollary}
\begin{proof}
由互反,
$\dim\End_G(\Ind W)=\dim\Hom_H(W,\Res_H^G\Ind W)$. 再用Mackey分解与另一方向的互反, 右侧为所有双陪集对应的上述Hom维数之和. 单位双陪集可取 $t=e$, 对应 $\End_H(W)$, 维数为1. 其余项均为非负整数, 所以总维数为1恰当且仅当它们都为0. 完全可约表示的交织代数维数为重数平方和, 等于1恰对应不可约.
\end{proof}
当 $H\trianglelefteq G$, 每个 $L_t=H$, 所以判据简化为 $W$ 与各非平凡商陪集给出的共轭表示互不等价. 若某个共轭表示与 $W$ 等价, 诱导中就存在额外交织算子, 因而会分裂.

\originalsection{一般二面体群}
令 $D_{2n}=\langle r,s:r^n=s^2=e,\ srs=r^{-1}\rangle$, $n\ge3$, 阶为 $2n$. 旋转子群 $H=\langle r\rangle\cong C_n$ 正规且指数为2. 令 $\zeta=\exp(2\pi i/n)$, $L_j(r)=\zeta^j$, 并记 $E_j=\Ind_H^{D_{2n}}L_j$.

选陪集代表 $e,s$. 在相应两份空间的基中,
\[
 \rho_{E_j}(r)=\begin{pmatrix}\zeta^j&0\\0&\zeta^{-j}\end{pmatrix},\qquad
 \rho_{E_j}(s)=\begin{pmatrix}0&1\\1&0\end{pmatrix}.
\]
限制到 $H$ 是 $L_j\oplus L_{-j}$; $s$ 交换两条特征直线. 当 $j\not\equiv-j\pmod n$ 时, 二面体作用没有不变直线, 因为任何 $r$ 不变直线都只能是两条特征直线之一, 而它们被 $s$ 交换. 因而 $E_j$ 不可约. 这也是上述正规子群判据的具体情形.

若 $j=0$, $r$ 作用恒等, $s$ 的 $\pm1$ 特征线给出两个一维表示. 若 $n$ 为偶数且 $j=n/2$, $r$ 作用为 $-I$, 同样分成 $s$ 取 $\pm1$ 的两个一维表示. 其他 $E_j$ 与 $E_{-j}$ 等价, 用交换两个基向量实现; 若两组 $\{j,-j\}$ 不同, 限制到 $H$ 的特征标不同, 故不等价.

\begin{theorem}
若 $n$ 为奇数, $D_{2n}$ 有两个一维表示, 以及 $E_j$, $1\le j\le(n-1)/2$. 若 $n$ 为偶数, 它有四个一维表示, 以及 $E_j$, $1\le j<n/2$. 这些列表穷尽全部复不可约表示. 二维表示的特征标为
\[
 \chi_{E_j}(r^a)=\zeta^{ja}+\zeta^{-ja}=2\cos(2\pi ja/n),\qquad
 \chi_{E_j}(r^as)=0.
\]
\end{theorem}
\begin{proof}
一维表示满足 $r\mapsto a=a^{-1}$, $s\mapsto b=\pm1$, 且 $a^n=1$. 奇数 $n$ 只能 $a=1$, 偶数 $n$ 可取 $a=\pm1$, 所以分别有二个或四个. 上面已构造并验证二维不可约及互不等价. 奇数时维数平方和为 $2+4(n-1)/2=2n$; 偶数时为 $4+4(n/2-1)=2n$. 所以分类完整. 特征标由所列矩阵直接取迹得到.
\end{proof}
取 $n=4$ 即回到第五章的 $D_8$; 取 $n=3$ 得 $D_6\cong S_3$. 一般分类把这些小群的表统一起来.

\originalsection{限制与分支计算}
\begin{example}
设 $m\mid n$ 且 $m\ge3$, $D_{2m}=\langle r^{n/m},s\rangle\le D_{2n}$. 求 $E_j$ 限制到这个子群的分解.
\end{example}
\begin{solution}
$u=r^{n/m}$ 的矩阵对角元为 $\exp(2\pi i j/m)$ 与其逆, $s$ 仍交换两个基向量. 因此若 $j\not\equiv-j\pmod m$, 限制就是 $D_{2m}$ 的二维表示 $E_{j\bmod m}$. 若 $j\equiv0\pmod m$, 它分成 $u$ 取1、$s$ 分别取 $\pm1$ 的两个一维表示. 若 $m$ 为偶数且 $j\equiv m/2\pmod m$, 则分成 $u$ 取 $-1$、$s$ 分别取 $\pm1$ 的两个一维表示. 这说明大群的不可约可以在子群上出现完整或分裂两种行为.
\end{solution}

\begin{exercise}
设 $H\trianglelefteq G$, $W$ 为 $H$ 不可约表示, 各商陪集给出的共轭表示互不等价. 证明 $\Res_H^G\Ind_H^G W$ 是这些共轭表示各出现一次的直和, 并求诱导表示的维数.
\end{exercise}
\begin{solution}
正规性给 $H\cap tHt^{-1}=H$. Mackey分解每项都为 $W_t$, 代表 $t$ 遍历 $G/H$. 假设保证它们两两不等价, 所以限制中各出现一次, 同时判据说明诱导不可约. 维数为 $[G:H]\dim W$.
\end{solution}
